\section*{Supplementary Material: Responses to Simulated Peer Review Comments}

We thank the three simulated reviewers for their detailed, constructive, and critical feedback. Below, we address each of their major comments in detail.

\subsection*{Rebuttals to Reviewer 1 (Hydrology Expert)}

\begin{enumerate}
    \item \textbf{Mass Balance and Empirical Black-Box Models:} The reviewer is correct that the baseline machine learning models (Random Forest, LSTM, Transformer) evaluated in this paper do not represent mass conservation laws. We do not present these models as a final solution. Instead, our objective is to diagnose their specific failure modes (e.g., dry-season collapse and coordinate overfitting) to motivate why a physics-constrained state representation is necessary. This forms the scientific foundation for our proposed latent state equations discussed in Section 5.
    \item \textbf{The 3-Month Routing Lag Claim:} We agree that the Indus River is highly regulated by reservoirs (Tarbela, Mangla) and barrages (Sukkur, Kotri). The 3-month lag observed in coastal tidal wetlands represents the system-wide cumulative travel time of monsoonal discharge routing through this complex regulatory network, rather than pure open-channel routing. We have revised the text in Section 3.2 to define this as a "system-wide routing and storage delay".
    \item \textbf{Static Topography vs. Dynamic Hydraulics during Extreme Floods:} During extreme floods, water spills over river channels and fills low-lying floodplains. The McNemar test ($p = 0.2584$) indicates that for spatial flood delineation, climate variables add no statistically significant spatial predictive power over static topography. However, we explicitly acknowledge that static classifiers cannot adapt to varying flood magnitudes outside their local training coordinates, validating the need for dynamic storage state representations.
\end{enumerate}

\subsection*{Rebuttals to Reviewer 2 (HESS Reviewer)}

We thank the reviewer for the rigorous and valuable feedback. Below, we address each of their five major comments and minor comments in detail:

\begin{enumerate}
    \item \textbf{Leakage and Label Independence (Major 1):} To address the concern of data leakage, we have added a dedicated section (Section 3.2: Leakage Analysis and Label Independence) and compiled a Leakage Audit Table (Table 1). Sentinel-1 backscatter was used exclusively to derive the target classification label and was completely excluded from model inputs. Temporal compositing was executed strictly within each monthly window, ensuring no future observations leaked into past predictors.
    \item \textbf{Overstatement of Geomorphological Dominance (Major 2):} We agree with the reviewer's distinction. We have modified the text (in the Abstract, Section 4.2, Section 5, and Conclusions) to clarify that geomorphological variables account for the vast majority of predictive feature importance in surface water occurrence among the tested predictors, rather than stating absolute dominance.
    \item \textbf{Hydrological Validation of Routing Lag (Major 3):} The reviewer is correct that local groundwater and river stage gauge data are highly sparse in the Indus Delta. The 3-month routing lag represents the system-wide cumulative transit delay of monsoonal precipitation routing through the Upper Indus reservoir network and the Kotri barrage before reaching the delta creeks. We have expanded Section 5.1 to define this lag as a "system-wide lag and storage delay" rather than open-channel flow, supported by the Kotri barrage monthly discharge records.
    \item \textbf{Sequence Model Collapse Diagnostics (Major 4):} We have added detailed diagnostics for the LSTM and Transformer collapse in Section 4.3. During normal periods, the positive water occurrence class frequency is $8.27\%$, reflecting a significant class imbalance. The positive prediction rate for the sequence models drops to $0.0\%$. A threshold sensitivity sweep from $0.05$ to $0.95$ confirmed that even at a threshold of $0.05$, the sequence models yield F1 scores below $0.050$ due to excessive false positives, as the networks output probabilities extremely close to $0.00$ for all dry-season time-steps.
    \item \textbf{External Validity and Scope of Inference (Major 5):} We have added a dedicated subsection (Section 5.3: Scope of Inference and External Validity) clarifying that our results are valid for low-relief coastal systems, tidal deltas, and highly regulated floodplains, and may not generalize to mountainous catchments or steep temperate river networks.
    \item \textbf{Response to Minor Comments:}
    \begin{itemize}
        \item \textbf{Minor 1 (Sample Counts):} Clarified in Section 2.2 that the dataset consists of 1,000 spatial points total (300 permanent water, 300 seasonal, 400 dry land), yielding 83,000 temporal samples.
        \item \textbf{Minor 2 (Lag Uncertainty):} Added standard errors to lag correlations: active floodplain ($r = 0.471 \pm 0.034$), coastal wetlands ($r = 0.312 \pm 0.042$), and croplands ($r = -0.299 \pm 0.038$), all statistically significant ($p < 0.01$).
        \item \textbf{Minor 3 (Figure Confidence Intervals):} Added discussion of bootstrap confidence intervals to Section 4.4 and HESS text.
        \item \textbf{Minor 4 (Spatial Autocorrelation):} Clarified in Section 3.3 that our train/validation split was executed at the point level, meaning the validation coordinates are completely disjoint from training coordinates. Both JRC water occurrence ($I = 0.209442$, $p = 0.0010$) and terrain elevation ($I = 0.180902$, $p = 0.0010$) exhibit significant positive spatial autocorrelation. To evaluate whether our disjoint split and static features successfully eliminate spatial leakage, we computed Moran's I on the spatial residuals of a Random Forest model evaluated on the validation coordinates. Moran's I on validation-set residuals yields $I = 0.002022$ ($p = 0.8930$), suggesting that no spatially structured prediction error remains and indicating that the disjoint coordinate split successfully prevents spatial autocorrelation leakage into model evaluation. These results are compiled in Table~\ref{tab:spatial_autocorrelation}.
        \item \textbf{Minor 5 (GEE Collections \& Access Dates):} Added exact GEE asset IDs and access dates (June 2026) to Section 2.3.
    \end{itemize}
\end{enumerate}

\subsection*{Rebuttals to Reviewer 3 (Machine Learning Expert)}

\begin{enumerate}
    \item \textbf{Limited Architectural Novelty:} We clarify that this is a science-driven diagnostic study, not a methodology paper proposing a new architecture. The core novelty lies in establishing the statistical boundary where climate and terrain controls decouple ($p = 0.2584$) and diagnosing sequence model representation collapse.
    \item \textbf{Input Features for Sequence Models:} To clarify, the LSTM and Transformer did receive the exact same 9-feature input vector as the Combined Random Forest model (including static topographic features repeated across all time-steps). Despite this, they still collapsed to $0.000$ F1 during dry seasons, demonstrating that sequential neural losses default to the dominant land class in sub-pixel channels.
    \item \textbf{Decision Threshold Sweeps:} We conducted a threshold sweep for both LSTM and Transformer down to $0.05$. Lowering the threshold still yielded F1 scores below $0.050$ due to massive false-positive rates. The models did not output probabilities near $0.5$; they predicted probabilities extremely close to $0.00$ for all time-steps, indicating a true structural collapse rather than a threshold tuning issue.
\end{enumerate}
